Knielend voor de kwade:
"Manlief, heb genade!”

De huishoudbeurs bleek leeg.
Vaders opwinding steeg.

Alles met geldzaken
wist hem hard te raken.

“Weer gesnoept uit de kas?”
Hoewel het nodig was,

kreeg ze op haar flikker.
De stok was niet dikker

dan zijn zuigende duim.
Vaders vuurmond vol schuim

schoot zijn geldmonoloog,
hoe ieder hem bedroog.

Zijn angsten en blamen
vuurden op haar samen

om ‘Moeder’ te raken.
“Mag ik het goedmaken?”

Want conform Gods wensen,
die de Kolossenzen

3 vers 18 melden,
moest voor Moeder gelden

onderdanig te zijn
aan haar ega’s welzijn.

Zijn spervuurstorm luwde
toen ze haar hoofd duwde

onder vaders schaamklep.
“Je weet wat ik lief heb!",

gaf hij haar als lintje.
Hij liet een klein windje

zijn woede uitrazen.
Echter bij het blazen

kwam er wat diarree
uit zijn poepgaatje mee.

Dieren ontweken hem.
Een bange kinderstem

verried zijn ijzigheid.
Een duistere kilheid,

die ineens ontremmend
Emmy luchtig zwemmend

onder water duwde.
En toen ze uitspuwde

haar bittere vloeistof,
onschuldig deed alsof

ze zich weer aanstelde.
Emmy opnieuw kwelde

met korte ademnood.
Zijn humeurigheid bood

geen grond voor connectie.
Daarom bezag Emmy

hem slechts als een huisvoogd,
die ze steeds had gedoogd.

Een buurkind zei een keer
“Ik zou niet graag, meneer,

met u getrouwd zijn, hoor!”
“Wees daar maar niet bang voor.

Dat zal nooit gebeuren!”
Ze staakte haar zeuren.

“Waarom koos je papa?”
vroeg Emmy haar mama.

“Hij kon ooit waarborgen
voor ons te gaan zorgen.

God leert ons dat liefde,
zelfs voor die ons griefde,

niet selectief mag zijn.
Eros, philia zijn schijn

en geenszins goddelijk.
Ze zijn voorwaardelijk.

Lees deze wijze les,
maar eens in Lukas 6.

Dus geef ik me over
en doe zand erover.”

Toch leek ma onzeker
of wel de gifbeker,

die ze ooit moest drinken,
haar niet zou verdrinken.

Ze sprong in het geloof
dat God haar heil toe schoof.

Later vond Kierkegaard
dit vast een pluimpje waard.

Terwijl ze hem afzoog
keek ze peilend omhoog

naar haar verweerde man.
De getemde tiran

tuurde naar het zogen
met ijsblauwe ogen.

Priemend keurde hij haar.
Door een trekhandgebaar,

met een knokige hand,
in moeders nek geplant,

nam hij weer de leiding
bij de vereffening.

Ingedeukte wangen
trilden van verlangen

onder zijn dunne baard.
Wit en golvend behaard,

zonder ooit te kammen,
toonden zijn inhammen.

Diepe rimpels kerfden
en zonnebrand verfden

zijn aangetrokken huid.
Vaders helm schoof eruit

om zaad te gaan knallen
uit zijn miniballen.

Daarop keerde doorbloed
de ingevallen snoet

zijn aanblik naar boven.
Jukbeenderen schoven

hun vlijmscherpe boegspriet
naar het plafond van riet.

Thomas van Aquino
zou de fellatio

als onnatuurlijk zien.
Vanuit Onan bezien,

niet eens interruptus,
want zelfs geen coïtus.

Hij ontving moeders blik,
maar leegde toch zijn pik.

“Verberg je verwijten
of het zal je spijten.

God heeft al besloten
en je schoot gesloten

als Sara, maar zonder
het conceptiewonder

net als bij mijn ex-vrouw.
Die had slechts winterkou.
Jij had ooit de lente.”
Waarbij hij inprentte

het verzonnen excuus.
“Al dit zaad was voor Truus”

Ook na zijn verlangen
bleef zijn wurggreep hangen.

Alsof vader klem zat
en pas weer de kracht had

zichzelf te ontspannen
na een rusttijdspanne.

“Niets zeggen hoor, Emmy!”,
zei ‘Moeder’ in frictie

met de realiteit.
De thuisvoogd zal geheid

het tekort ontdekken.
Het was slechts tijd rekken.

“Papa wordt dan zo boos.
Er is waarlijk niets loos

als hij de landheer smeekt.
De pacht ons niet opbreekt.

Maar hij wil dat niet doen
uit aanzien en fatsoen.

Laat zich toe-eigenen
als lijfeigene

als we niet ingrijpen.
Kun je dat begrijpen?

Hij is een ja-knikker.
Maar ook een driftkikker.

Kun jij hem niet vragen?
Dat kan ik niet wagen!"

Zo doorleefde Emmy
parentificatie.

“Hoeveel verdient hij dan?”
“Hij verdient als mijn man

wel duizend Rijkstaler!”
Voor een pachtbetaler

was dat meer dan genoeg.
Vandaar dat Emmy vroeg

naar wat hem had bedreigd.
“Het is niet wat hij krijgt!”

Slechts als komediant
smeedde Emmy een band.

Samen met haar thuisvoogd
werd ondeugend gepoogd

mensen te verwarren
of ze speels te sarren.

Ze improviseerden
of accentueerden

typerende rollen
om samen te dollen

rond naïeve mensen
met respect voor grenzen.

Emmy hoorde een keer
de abdis in de weer

bij de lekenzusters.
Ze gaf hen dan clusters

van telkens drie opties.
Koos men vervolgacties,

die er niet in pasten,
dan hoorde de gasten

haar felle: “Nee, nee, nee!”
En nam ze hen weer mee

in de drie vrijheden.
Wat de acteurs deden

was op de kloosterweg
te staan met quasi pech

gelijk warenwoekers
om kloosterbezoekers

te vragen om bijstand.
Wees men dat van de hand,

vanwege het beroep,
dan volgde de oproep

theatraal en te luid:
“Nee, nee, nee! Kiest u uit

lachen, gieren, brullen.”
Vader mocht onthullen

met een brede grimas
dat alles een spel was.

Op een zekere dag
toen pa op de loer lag

tesamen met Emmy
voor een plaisanterie

met een abdij logé
observeerden de twee

een verduisterde koets.
Ze begonnen hun poets
tot er uit het rijtuig
“Wat een respectloos tuig!”

kwam als gekrenkte schreeuw.
In het hol van de leeuw

brulde de abdis zelf.
De thuisvoogd ging vanzelf

het volksvermaak staken
en excuses maken.

Tot ze zei: “Nee, nee, nee!”
en toen lachte ze mee.

Deed ze dat met kiespijn?
En kwam Emmy’s welzijn

ooit hiermee in gevaar?
U wordt het vast gewaar!
